% !TEX program = xelatex
% 공통수학2 · Ⅲ. 함수와 그래프 · 1. 함수 · 02 합성함수 · 1/1차시
% 학생용 활동지 (답안 없음)
\documentclass[11pt,a4paper]{article}
\usepackage{kotex}
\usepackage{iftex}
\ifPDFTeX\else
  \setmainhangulfont{Noto Serif CJK KR}
  \setsanshangulfont{Noto Sans CJK KR}
\fi
\usepackage[margin=2.1cm]{geometry}
\usepackage{parskip}
\usepackage{enumitem}
\usepackage{array}
\usepackage{tabularx}
\usepackage{xcolor}
\usepackage{amssymb}
\usepackage{tikz}
\usepackage[most]{tcolorbox}
\usepackage{fancyhdr}
\usepackage{titlesec}

\definecolor{goldc}{HTML}{B07C22}
\definecolor{goldstrong}{HTML}{8A5F16}
\definecolor{indigoc}{HTML}{2B3B57}
\definecolor{indigosoft}{HTML}{6E82A6}
\definecolor{linec}{HTML}{D6DACB}
\definecolor{surface2}{HTML}{E4E8DC}
\definecolor{writeline}{HTML}{C7CCBA}
\definecolor{inksoft}{HTML}{52604F}

\titleformat{\section}{\normalfont\sffamily\bfseries\large\color{indigoc}}{\thesection}{0.6em}{}
\titlespacing{\section}{0pt}{14pt}{6pt}
\newtcolorbox{goalbox}{enhanced, breakable, colback=white, colframe=goldc, boxrule=0.5pt, leftrule=3pt, arc=2pt, left=10pt, right=10pt, top=8pt, bottom=8pt}
\newtcolorbox{sheetbox}{enhanced, breakable, colback=white, colframe=linec, boxrule=0.6pt, arc=3pt, left=11pt, right=11pt, top=9pt, bottom=9pt}
\newcommand{\writespace}[1]{%
  \par\nobreak\vspace{3pt}
  \foreach \n in {1,...,#1} {%
    \noindent\textcolor{writeline}{\rule{\linewidth}{0.4pt}}\par\vspace{0.62cm}%
  }
}
\newcommand{\fillblank}[1]{\makebox[#1]{\hrulefill}}
\pagestyle{fancy}\fancyhf{}\renewcommand{\headrulewidth}{0pt}
\fancyfoot[L]{\footnotesize\color{inksoft} 공통수학2 · 2026학년도 2학기 · 지평선고등학교}
\fancyfoot[R]{\footnotesize\color{inksoft} 다음 시간 --- 03 역함수 1/2}

\newcommand{\chainmap}[1]{%
  \begin{tikzpicture}[scale=0.62, baseline]
    \draw[thick, indigosoft] (0,0) ellipse (0.75 and 1.3);
    \draw[thick, indigosoft] (2.2,0) ellipse (0.75 and 1.3);
    \draw[thick, indigosoft] (4.4,0) ellipse (0.75 and 1.3);
    \node[above] at (0,1.5) {\footnotesize $X$};
    \node[above] at (2.2,1.5) {\footnotesize $Y$};
    \node[above] at (4.4,1.5) {\footnotesize $Z$};
    \node[below] at (1.1,-1.6) {\footnotesize $f$};
    \node[below] at (3.3,-1.6) {\footnotesize $g$};
    #1
  \end{tikzpicture}}

\begin{document}

\noindent
\begin{minipage}[b]{0.62\linewidth}
  {\sffamily\footnotesize\color{indigosoft} 공통수학2 · Ⅲ. 함수와 그래프 · 1. 함수}\\[2pt]
  {\sffamily\bfseries\Large 02 합성함수 \textperiodcentered\ 35차시}
\end{minipage}\hfill
\begin{minipage}[b]{0.36\linewidth}
  \raggedleft\small\color{inksoft}
  반\ \fillblank{1.6cm} \quad 번호\ \fillblank{1.6cm}\\[4pt]
  이름\ \fillblank{3.6cm}
\end{minipage}

\vspace{4pt}
\noindent{\color{goldc}\rule{\linewidth}{2.2pt}}
\vspace{10pt}

\begin{goalbox}
\textbf{오늘의 목표}\\
합성함수의 뜻을 알고, 두 함수의 합성함수를 구하며 순서를 바꾸면 결과가 달라짐을 확인할 수 있다.
\vspace{4pt}
\small\color{inksoft}
$\square$ 자신 있음 \quad $\square$ 조금 헷갈림 \quad $\square$ 처음 봄 \hfill (수업 전 나의 상태에 표시)
\end{goalbox}

\section{생각 열기 --- 순서를 바꾸면?}

\begin{sheetbox}
다음 두 행동의 순서를 비교해 보자.

\begin{enumerate}[leftmargin=1.4em, itemsep=2pt]
  \item 양말을 신는다 $\to$ 그 위에 신발을 신는다.
  \item 신발을 신는다 $\to$ 그 위에 양말을 신는다.
\end{enumerate}

두 순서 중 어느 쪽이 `말이 되는' 순서인가? 이것을 함수를 두 번 잇는 상황과 연결하면 어떤 예감이 드는지 써 보자.
\writespace{3}
\end{sheetbox}

\section{개념 정리 --- 합성함수}

\begin{sheetbox}
두 함수 $f:X\to Y$, $g:Y\to Z$에 대하여, $X$의 각 원소 $x$에 $g(f(x))$를 대응시켜 만든 새로운 함수를 $f$와 $g$의 \fillblank{2.6cm}라 하고, 기호로 \fillblank{2cm}와 같이 나타낸다.

\vspace{6pt}
\[(g\circ f)(x) = g(\ \fillblank{1.6cm}\ )\]

\vspace{4pt}
\noindent
\begin{center}
\chainmap{}
\end{center}

\vspace{6pt}
계산할 때는 \fillblank{1.6cm} 함수를 먼저 적용하고, 그 결과에 \fillblank{1.6cm} 함수를 적용한다.

\vspace{8pt}
\textbf{예제} --- $f(x)=x+1$, $g(x)=2x$일 때\\
$(g\circ f)(3) = g(f(3)) = g(\fillblank{1.2cm}) = \fillblank{1.2cm}$
\end{sheetbox}

\section{연습 문제}

\begin{sheetbox}
\textbf{1.} $f(x)=2x+1$, $g(x)=x^2$일 때, $(g\circ f)(2)$와 $(f\circ g)(2)$를 각각 구하고, 두 값을 비교하시오.
\writespace{4}

\vspace{4pt}
\textbf{2.} $f(x)=x-3$, $g(x)=2x+1$일 때, $(g\circ f)(x)$와 $(f\circ g)(x)$를 각각 구하고, 서로 같은지 판별하시오.
\writespace{4}

\vspace{4pt}
\textbf{3.} $f(x)=x+2$, $g(x)=3x$, $h(x)=x^2$일 때, $((h\circ g)\circ f)(x)$와 $(h\circ(g\circ f))(x)$를 각각 구하여 두 식이 같은지 확인하시오.
\writespace{5}

\vspace{4pt}
\textbf{4.} $X=\{1,2,3\}$, $Y=\{2,4,6\}$, $Z=\{4,16,36\}$이고 $f:X\to Y$, $f(x)=2x$, $g:Y\to Z$, $g(y)=y^2$일 때, $(g\circ f)(3)$을 구하고, $g\circ f$가 정의될 수 있는 이유를 한 문장으로 설명하시오.
\writespace{4}
\end{sheetbox}

\section{사고의 가시화 --- Connect · Extend · Challenge}

\noindent{\footnotesize\color{inksoft} 오늘 배운 `합성함수'를 떠올리며 아래 세 칸을 채워 보자.}\\[6pt]
\begin{tabularx}{\linewidth}{@{}XXX@{}}
\textbf{\footnotesize\color{indigoc} Connect --- 무엇과 연결되나?} &
\textbf{\footnotesize\color{indigoc} Extend --- 어디까지 확장됐나?} &
\textbf{\footnotesize\color{indigoc} Challenge --- 아직 궁금한 것은?} \\[4pt]
\end{tabularx}
\noindent
\begin{minipage}[t]{0.32\linewidth}\writespace{3}\end{minipage}\hfill
\begin{minipage}[t]{0.32\linewidth}\writespace{3}\end{minipage}\hfill
\begin{minipage}[t]{0.32\linewidth}\writespace{3}\end{minipage}

\section{마무리 성찰}

\begin{sheetbox}
\begin{minipage}[t]{0.48\linewidth}
{\footnotesize\color{inksoft} $g\circ f$와 $f\circ g$가 다른 이유를 한 문장으로}
\writespace{2}
\end{minipage}\hfill
\begin{minipage}[t]{0.48\linewidth}
{\footnotesize\color{inksoft} 감사 일기 / 유념 한 줄}
\writespace{2}
\end{minipage}
\end{sheetbox}

\end{document}
